\section{Analisis Performa Portofolio Akhir: Perbandingan GT-VQE vs \textit{Benchmark}}
\label{sec:4_10_performa_akhir}

Keunggulan strategis dari model hibrida \textit{Game Theory}-VQE dievaluasi melalui pengujian \textit{backtesting} yang merefleksikan kondisi pasar modal Indonesia yang dinamis.

\subsection{Penurunan Rumus Metrik Finansial: Imbal Hasil dan Risiko}
\label{subsec:4_10_1_rumus}

Performa portofolio diukur secara kuantitatif menggunakan tiga metrik utama: \textit{Cumulative Return}, \textit{Sharpe Ratio}, dan \textit{Maximum Drawdown} (MDD). Imbal hasil kumulatif $R_c$ dihitung melalui akumulasi pertumbuhan nilai aset dari waktu ke waktu, sementara \textit{Sharpe Ratio} ($S$) diturunkan dari rasio antara ekspektasi imbal hasil portofolio $\mu_p$ terhadap volatilitas portofolio $\sigma_p$, setelah dikurangi tingkat bunga bebas risiko $r_f$:
\begin{equation}
    S = \frac{\mu_p - r_f}{\sigma_p}
\end{equation}
Metrik MDD mengukur penurunan terbesar dari titik puncak ke titik terendah selama periode investasi, yang menjadi indikator ketangguhan strategi dalam menghadapi fase pasar yang terkontraksi.

\subsection{Analisis Sensitivitas terhadap Variasi Kondisi Pasar}
\label{subsec:4_10_2_sensitivitas}

Analisis sensitivitas dilakukan untuk mengamati stabilitas performa algoritma dalam berbagai rezim pasar (\textit{bullish}, \textit{bearish}, dan \textit{sideways}). Hasil pengujian menunjukkan bahwa strategi GT-VQE memiliki ketahanan yang lebih baik dibandingkan strategi \textit{Equal Weight} (EW) saat terjadi guncangan pasar. Hal ini disebabkan oleh kemampuan algoritma untuk secara dinamis merestrukturisasi bobot aset berdasarkan perubahan profil risiko-return yang dipetakan ke dalam Hamiltonian. Penyesuaian parameter \textit{risk aversion} endogen memungkinkan portofolio untuk melakukan mitigasi kerugian secara proaktif, sehingga menjaga volatilitas kurva ekuitas tetap berada dalam batas toleransi investasi.

\subsection{Simulasi Pertumbuhan Modal Investasi}
\label{subsec:4_10_3_numerik}

Simulasi pertumbuhan modal investasi awal disajikan secara visual pada Gambar \ref{fig:backtest_curve_vqe_gt} di Lampiran \ref{appendix:additional_results}. Berdasarkan laporan kinerja pada Tabel \ref{tab:metrics_unified}, strategi GT-VQE mencatatkan total imbal hasil kumulatif sebesar 51,34\%, mengungguli IHSG yang hanya tumbuh sebesar 19,13\% pada periode yang sama. Keunggulan kuantitatif ini membuktikan bahwa integrasi kecerdasan strategis \textit{Game Theory} ke dalam sistem komputasi kuantum memberikan nilai tambah yang signifikan dalam meningkatkan efisiensi alokasi modal dan memaksimalkan \textit{risk-adjusted return} bagi investor profesional.
